\documentclass[a4paper, serif, 10pt]{moderncv}

%% moderncv
% style and color
\moderncvstyle{banking}
\moderncvcolor{black}

%% page layout/margins
\usepackage[top=2.5cm, bottom=2.5cm, left=1.5cm, right=1.5cm]{geometry}

\nopagenumbers{}

%% Babel config for Spanish language
% \usepackage[spanish]{babel} has some conflicting issues with moderncv
% so we have to use enable the following options to make it work
%
% ref:
% - https://tex.stackexchange.com/a/140161/36007
\usepackage[spanish,es-lcroman]{babel}

%% fontspec
\usepackage{fontspec}

\IfFontExistsTF{Linux Libertine}{
  \setmainfont[Ligatures={TeX, Common}, Numbers=OldStyle]{Linux Libertine}
}{}
\IfFontExistsTF{Linux Libertine O}{
  \setmainfont[Ligatures={TeX, Common}, Numbers=OldStyle]{Linux Libertine O}
}{}

%% CTeX
% CJK support, used to show CJK characters in the resume
%
% - fontset=none: disable builtin fontset but instead set the CJK font manually
% - heading=false: disable ctex heading
% - punct=kaiming: use kaiming punctuations styles for CJK
% - scheme=plain: use plain scheme, do not override \normalsize font size
% - space=auto: space settings for CJK characters
%
% ref:
% - http://ctan.mirrorcatalogs.com/language/chinese/ctex/ctex.pdf
\usepackage[UTF8, heading=false, punct=kaiming, scheme=plain, space=auto]{ctex}

\IfFontExistsTF{Noto Serif CJK SC}{
  \setCJKmainfont{Noto Serif CJK SC}
}{}
\IfFontExistsTF{Noto Sans CJK SC}{
  \setCJKsansfont{Noto Sans CJK SC}
}{}

%% line spacing
\usepackage{setspace}
\setstretch{1.125}

%% URL styling - use normal text instead of monospace
\urlstyle{same}

% Auto-underline all links
% Defer until \begin{document} so hyperref is already loaded
\AtBeginDocument{%
  \let\oldhref\href
  \renewcommand{\href}[2]{\oldhref{#1}{\underline{#2}}}%
}

%% Patch \httplink and \httpslink to handle full URLs
% The original moderncv macros always prepend http:// or https:// to URLs.
% This causes issues when the URL already has a protocol (e.g., https://example.com)
% resulting in malformed URLs like https://https://example.com.
% This patch checks if the URL already contains "://" and uses it directly if so.
% Note: We use \str_set:Nx (x-expansion) to expand macros like \@homepage first.
\ExplSyntaxOn
\str_new:N \g_yamlresume_url_str

\RenewDocumentCommand{\httpslink}{O{}m}{
  \str_set:Nx \g_yamlresume_url_str {#2}
  \str_if_in:NnTF \g_yamlresume_url_str {://}
    { \url{#2} }
    { \url{https://#2} }
}

\RenewDocumentCommand{\httplink}{O{}m}{
  \str_set:Nx \g_yamlresume_url_str {#2}
  \str_if_in:NnTF \g_yamlresume_url_str {://}
    { \url{#2} }
    { \url{http://#2} }
}
\ExplSyntaxOff

\name{Lucía Fernández}{}
\title{Diseñadora UX/UI senior, especializada en investigación y diseño de servicios digitales}
\phone[mobile]{+34 612 345 678}
\email{lucia@fernandezux.com}
\homepage{https://www.fernandezux.com}

\address{Calle Gran Vía 28, 2º, Madrid, Comunidad de Madrid, España, 28013}{}{}

\extrainfo{{\small \faLinkedin}\ \href{https://www.linkedin.com/in/luciafernandezux}{@luciafernandezux} {} {} {} • {} {} {}
{\small \faBehance}\ \href{https://www.behance.net/luciafernandez}{@luciafernandez}}

\begin{document}

\maketitle

\section{Información básica}

\cvline{}{Diseñadora de producto con más de 7 años de experiencia en proyectos digitales para banca, salud y retail.
Combino investigación cualitativa, análisis de datos y prototipado para crear experiencias usables y accesibles.
%
\begin{itemize}
\item Lidero el diseño de extremo a extremo: research, ideación, prototipado, testing y traspaso a desarrollo.
\item Defiendo la accesibilidad (WCAG 2.2) y el diseño inclusivo como parte esencial del proceso.
\item Colaboro con equipos ágiles, producto y negocio para alinear objetivos y prioridades.
\end{itemize}}
\section{Educación}

\cventry{sept 2016–jul 2018}
        {Maestría, Diseño de Experiencia de Usuario}
        {Universidad Politécnica de Madrid}
        {\url{https://www.upm.es/}}
        {}
        {\begin{itemize}
\item Proyecto final premiado: rediseño de una plataforma de teleasistencia para mayores.
\item Colaboración con el Laboratorio de Accesibilidad para auditar apps públicas.
\end{itemize}
\textbf{Cursos}: Métodos de investigación con usuarios, Diseño de interacción, Arquitectura de la información, Evaluación de usabilidad, Visualización de datos}

\cventry{sept 2011–jun 2015}
        {Licenciatura, Ingeniería de Diseño Industrial}
        {Universidad Carlos III de Madrid}
        {\url{https://www.uc3m.es/}}
        {}
        {\textbf{Cursos}: Ergonomía y factores humanos, Diseño centrado en el usuario, Tecnologías multimedia}
\section{Experiencia laboral}

\cventry{mar 2021 hasta la fecha}
        {Senior UX Designer}
        {Nubika Digital}
        {\url{https://www.nubikadigital.com}}
        {}
        {\begin{itemize}
\item Lidero el área de UX dentro de una squad de banca digital, colaborando con PMs, analistas y desarrolladores.
\item Rediseño integral de la app móvil de un banco europeo, logrando un aumento del 34 \% en finalización de onboarding.
\item Creo y mantengo el design system junto a otros diseñadores, con más de 80 componentes en Figma.
\item Planifico y facilito sesiones de co-creación y testeos con personas usuarias en distintos países.
\end{itemize}
\textbf{Palabras clave}: Design Systems, Banca digital, Figma, Agile, Accessibility}

\cventry{jun 2018–feb 2021}
        {UX/UI Designer}
        {Pixel House Studio}
        {\url{https://www.pixelhousestudio.com}}
        {}
        {\begin{itemize}
\item Diseñé la experiencia de marca y producto para más de 20 clientes de los sectores salud, turismo y comercio electrónico.
\item Implementé procesos de investigación dentro de la agencia: entrevistas, encuestas y pruebas de usabilidad.
\item Trabajé junto a desarrolladores front-end para asegurar la viabilidad técnica de las propuestas.
\item Ganamos dos premios Lovie por proyectos de e-commerce accesibles.
\end{itemize}}
\section{Idiomas}

\cvline{Español}{Competencia nativa o bilingüe \hfill \textbf{Palabras clave}: Lengua materna}
\cvline{Inglés}{Competencia profesional plena \hfill \textbf{Palabras clave}: TOEFL 105, Inglés técnico}
\cvline{Francés}{Competencia limitada de trabajo}
\section{Competencias}

\cvline{Research}{Experto \hfill \textbf{Palabras clave}: Entrevistas con usuarios, Encuestas, Pruebas de usabilidad, Card sorting, Customer journeys}
\cvline{Prototipado y diseño}{Experto \hfill \textbf{Palabras clave}: Figma, Sketch, Prototipado interactivo, Wireframes, Design systems}
\cvline{Arquitectura de la información}{Avanzado \hfill \textbf{Palabras clave}: Sitemaps, Flujos de usuario, Navegación, Taxonomía}
\cvline{Accesibilidad}{Avanzado \hfill \textbf{Palabras clave}: WCAG 2.2, Diseño inclusivo, ARIA, Auditorías de accesibilidad}
\section{Premios}

\cventry{nov 2022}
        {Asociación de Diseño de Interacción (IxDA España)}
        {Premio Nacional de Diseño de Experiencia}
        {}
        {}
        {Reconocimiento al mejor proyecto de diseño de servicios digitales por la app de salud "Bienestar360".}
\section{Certificados}

\cventry{abr 2020}
        {Coursera}
        {Google UX Design Professional Certificate}
        {\url{https://www.coursera.org/professional-certificates/google-ux-design}}
        {}
        {}

\cventry{oct 2019}
        {Human Factors International}
        {Certified Usability Analyst}
        {\url{https://www.humanfactors.com}}
        {}
        {}
\section{Publicaciones}

\cventry{may 2023}
        {El impacto del design system en la productividad de equipos ágiles}
        {Revista Iberoamericana de Interacción Humano-Computadora}
        {\url{https://www.revistaiahc.org/estudio-design-system}}
        {}
        {Estudio de caso sobre la implantación de un design system en una entidad financiera y su efecto en la velocidad de desarrollo y la consistencia del producto.}
\section{Proyectos}

\cventry{ene 2022–dic 2022}
        {Rediseño completo de una plataforma de telemedicina para mejorar la adherencia al tratamiento.}
        {Bienestar360}
        {\url{https://www.bienestar360.example}}
        {}
        {\begin{itemize}
\item Dirigí la investigación con pacientes y profesionales de la salud para identificar puntos de fricción.
\item Testeé prototipos en alta fidelidad con 30 usuarios, iterando el flujo de gestión de citas y medicación.
\item El rediseño redujo el abandono de la plataforma un 28 \% y las consultas de soporte un 47 \%.
\end{itemize}
\textbf{Palabras clave}: Salud digital, Telemedicina, User research, Prototipado}

\cventry{mar 2021–jul 2023}
        {Sistema de diseño de código abierto para aplicaciones de banca y seguros.}
        {Design System "Atomika"}
        {\url{https://www.designsystem.atomika.example}}
        {}
        {\begin{itemize}
\item Definición de tokens, componentes y documentación para un sistema multiplataforma.
\item Incluye pautas de accesibilidad y ejemplos de código para React Native.
\item Adoptado por 5 equipos de producto, reduciendo el tiempo de diseño de nuevas funcionalidades.
\end{itemize}
\textbf{Palabras clave}: Design system, Tokens, React, Documentación}
\section{Intereses}

\cvline{Fotografía urbana}{Street photography, Composición, Luz natural}
\cvline{Senderismo}{Rutas de montaña, Naturaleza, Bienestar}
\section{Voluntariado}

\cventry{ene 2020–mar 2023}
        {Coordinadora de voluntarios}
        {UXPA España}
        {\url{https://www.uxpa.es}}
        {}
        {\begin{itemize}
\item Organicé encuentros y talleres sobre metodologías UX para la comunidad local.
\item Fui mentora en el programa de jóvenes diseñadores, asesorando a personas en sus primeros proyectos profesionales.
\end{itemize}}

\end{document}